% Options for packages loaded elsewhere
% Options for packages loaded elsewhere
\PassOptionsToPackage{unicode}{hyperref}
\PassOptionsToPackage{hyphens}{url}
\PassOptionsToPackage{dvipsnames,svgnames,x11names}{xcolor}
%
\documentclass[
  11pt,
  letterpaper,
  DIV=11,
  numbers=noendperiod]{scrartcl}
\usepackage{xcolor}
\usepackage[margin=2.5cm]{geometry}
\usepackage{amsmath,amssymb}
\setcounter{secnumdepth}{5}
\usepackage{iftex}
\ifPDFTeX
  \usepackage[T1]{fontenc}
  \usepackage[utf8]{inputenc}
  \usepackage{textcomp} % provide euro and other symbols
\else % if luatex or xetex
  \usepackage{unicode-math} % this also loads fontspec
  \defaultfontfeatures{Scale=MatchLowercase}
  \defaultfontfeatures[\rmfamily]{Ligatures=TeX,Scale=1}
\fi
\usepackage{lmodern}
\ifPDFTeX\else
  % xetex/luatex font selection
\fi
% Use upquote if available, for straight quotes in verbatim environments
\IfFileExists{upquote.sty}{\usepackage{upquote}}{}
\IfFileExists{microtype.sty}{% use microtype if available
  \usepackage[]{microtype}
  \UseMicrotypeSet[protrusion]{basicmath} % disable protrusion for tt fonts
}{}
\usepackage{setspace}
\makeatletter
\@ifundefined{KOMAClassName}{% if non-KOMA class
  \IfFileExists{parskip.sty}{%
    \usepackage{parskip}
  }{% else
    \setlength{\parindent}{0pt}
    \setlength{\parskip}{6pt plus 2pt minus 1pt}}
}{% if KOMA class
  \KOMAoptions{parskip=half}}
\makeatother
% Make \paragraph and \subparagraph free-standing
\makeatletter
\ifx\paragraph\undefined\else
  \let\oldparagraph\paragraph
  \renewcommand{\paragraph}{
    \@ifstar
      \xxxParagraphStar
      \xxxParagraphNoStar
  }
  \newcommand{\xxxParagraphStar}[1]{\oldparagraph*{#1}\mbox{}}
  \newcommand{\xxxParagraphNoStar}[1]{\oldparagraph{#1}\mbox{}}
\fi
\ifx\subparagraph\undefined\else
  \let\oldsubparagraph\subparagraph
  \renewcommand{\subparagraph}{
    \@ifstar
      \xxxSubParagraphStar
      \xxxSubParagraphNoStar
  }
  \newcommand{\xxxSubParagraphStar}[1]{\oldsubparagraph*{#1}\mbox{}}
  \newcommand{\xxxSubParagraphNoStar}[1]{\oldsubparagraph{#1}\mbox{}}
\fi
\makeatother


\usepackage{longtable,booktabs,array}
\newcounter{none} % for unnumbered tables
\usepackage{calc} % for calculating minipage widths
% Correct order of tables after \paragraph or \subparagraph
\usepackage{etoolbox}
\makeatletter
\patchcmd\longtable{\par}{\if@noskipsec\mbox{}\fi\par}{}{}
\makeatother
% Allow footnotes in longtable head/foot
\IfFileExists{footnotehyper.sty}{\usepackage{footnotehyper}}{\usepackage{footnote}}
\makesavenoteenv{longtable}
\usepackage{graphicx}
\makeatletter
\newsavebox\pandoc@box
\newcommand*\pandocbounded[1]{% scales image to fit in text height/width
  \sbox\pandoc@box{#1}%
  \Gscale@div\@tempa{\textheight}{\dimexpr\ht\pandoc@box+\dp\pandoc@box\relax}%
  \Gscale@div\@tempb{\linewidth}{\wd\pandoc@box}%
  \ifdim\@tempb\p@<\@tempa\p@\let\@tempa\@tempb\fi% select the smaller of both
  \ifdim\@tempa\p@<\p@\scalebox{\@tempa}{\usebox\pandoc@box}%
  \else\usebox{\pandoc@box}%
  \fi%
}
% Set default figure placement to htbp
\def\fps@figure{htbp}
\makeatother





\setlength{\emergencystretch}{3em} % prevent overfull lines

\providecommand{\tightlist}{%
  \setlength{\itemsep}{0pt}\setlength{\parskip}{0pt}}



 


\KOMAoption{captions}{tableheading}
\makeatletter
\@ifpackageloaded{caption}{}{\usepackage{caption}}
\AtBeginDocument{%
\ifdefined\contentsname
  \renewcommand*\contentsname{Table of contents}
\else
  \newcommand\contentsname{Table of contents}
\fi
\ifdefined\listfigurename
  \renewcommand*\listfigurename{List of Figures}
\else
  \newcommand\listfigurename{List of Figures}
\fi
\ifdefined\listtablename
  \renewcommand*\listtablename{List of Tables}
\else
  \newcommand\listtablename{List of Tables}
\fi
\ifdefined\figurename
  \renewcommand*\figurename{Figure}
\else
  \newcommand\figurename{Figure}
\fi
\ifdefined\tablename
  \renewcommand*\tablename{Table}
\else
  \newcommand\tablename{Table}
\fi
}
\@ifpackageloaded{float}{}{\usepackage{float}}
\floatstyle{ruled}
\@ifundefined{c@chapter}{\newfloat{codelisting}{h}{lop}}{\newfloat{codelisting}{h}{lop}[chapter]}
\floatname{codelisting}{Listing}
\newcommand*\listoflistings{\listof{codelisting}{List of Listings}}
\makeatother
\makeatletter
\makeatother
\makeatletter
\@ifpackageloaded{caption}{}{\usepackage{caption}}
\@ifpackageloaded{subcaption}{}{\usepackage{subcaption}}
\makeatother
\usepackage{bookmark}
\IfFileExists{xurl.sty}{\usepackage{xurl}}{} % add URL line breaks if available
\urlstyle{same}
\hypersetup{
  pdftitle={Análise da Distribuição de Profissionais de Dados no Brasil},
  pdfauthor={Eliza Freitas; Pedro Greco; Satie; Miguel Sant'ana},
  colorlinks=true,
  linkcolor={blue},
  filecolor={Maroon},
  citecolor={Blue},
  urlcolor={Blue},
  pdfcreator={LaTeX via pandoc}}


\title{Análise da Distribuição de Profissionais de Dados no Brasil}
\usepackage{etoolbox}
\makeatletter
\providecommand{\subtitle}[1]{% add subtitle to \maketitle
  \apptocmd{\@title}{\par {\large #1 \par}}{}{}
}
\makeatother
\subtitle{Trabalho Final - Recuperação da Informação}
\author{Eliza Freitas \and Pedro Greco \and Satie \and Miguel Sant'ana}
\date{}
\begin{document}
\maketitle

\renewcommand*\contentsname{Table of contents}
{
\setcounter{tocdepth}{3}
\tableofcontents
}

\setstretch{1.5}
\section{Identificação do Grupo}\label{identificauxe7uxe3o-do-grupo}

\textbf{Disciplina:} Ciência de Dados

\textbf{Curso:} Bacharelado em Ciência da Computação

\textbf{Professor:} Weslley Rodrigues

\textbf{Data de Entrega:} 22 de junho de 2026

\subsection{Integrantes}\label{integrantes}

{\def\LTcaptype{none} % do not increment counter
\begin{longtable}[]{@{}lll@{}}
\toprule\noalign{}
Nome & Matrícula & Usuário GitHub \\
\midrule\noalign{}
\endhead
\bottomrule\noalign{}
\endlastfoot
Eliza Freitas & 22302645 & Elizafc \\
Pedro Greco & 22308774 & GoGreco \\
Satie & 22301624 & sati-e \\
Miguel Sant'ana & 22304022 & Miguel-Sant77 \\
\end{longtable}
}

\begin{center}\rule{0.5\linewidth}{0.5pt}\end{center}

\section{Pergunta de Pesquisa}\label{pergunta-de-pesquisa}

\subsection{Problema de Pesquisa}\label{problema-de-pesquisa}

O mercado de dados no Brasil tem crescido significativamente nos últimos
anos, com um aumento exponencial na demanda por profissionais
qualificados. No entanto, essa distribuição não ocorre de forma
homogênea pelo território nacional, sendo influenciada por diversos
fatores econômicos e sociais.

\subsection{Pergunta Central}\label{pergunta-central}

\textbf{``Como a distribuição geográfica de profissionais de dados no
Brasil se relaciona com indicadores socioeconômicos como PIB, renda
média, IDH e custo de vida por estado?''}

\subsection{Objetivos Específicos}\label{objetivos-especuxedficos}

\begin{enumerate}
\def\labelenumi{\arabic{enumi}.}
\tightlist
\item
  Consolidar os dados do State of Data Brazil das edições de 2021 a 2024
\item
  Coletar indicadores socioeconômicos oficiais via API do IBGE
\item
  Integrar as duas bases de dados utilizando UF como chave
\item
  Aplicar modelo de Recuperação da Informação (TF-IDF) para busca de
  estados similares
\item
  Identificar correlações entre a concentração de profissionais e
  variáveis socioeconômicas
\item
  Gerar visualizações que auxiliem na interpretação dos resultados
\end{enumerate}

\subsection{Justificativa}\label{justificativa}

Compreender a distribuição geográfica dos profissionais de dados é
fundamental para:

\begin{itemize}
\tightlist
\item
  \textbf{Políticas Públicas:} Identificar regiões que necessitam de
  incentivos para formação de profissionais
\item
  \textbf{Empresas:} Auxiliar na decisão de onde instalar novos
  escritórios ou polos tecnológicos
\item
  \textbf{Profissionais:} Orientar decisões de carreira e realocação
  geográfica
\item
  \textbf{Academia:} Entender os fatores que impulsionam o mercado de
  dados no Brasil
\end{itemize}

\begin{center}\rule{0.5\linewidth}{0.5pt}\end{center}

\section{Descrição das Bases
Utilizadas}\label{descriuxe7uxe3o-das-bases-utilizadas}

\subsection{Base 1: State of Data Brazil
(2021-2024)}\label{base-1-state-of-data-brazil-2021-2024}

\subsubsection{Sobre a Pesquisa}\label{sobre-a-pesquisa}

O State of Data Brazil é uma pesquisa anual realizada pela Data Hackers,
que mapeia o perfil, comportamento e percepções dos profissionais que
trabalham com dados no Brasil. É a maior pesquisa do gênero no país e
serve como referência para o mercado de dados.

\subsubsection{Características da
Base}\label{caracteruxedsticas-da-base}

{\def\LTcaptype{none} % do not increment counter
\begin{longtable}[]{@{}
  >{\raggedright\arraybackslash}p{(\linewidth - 2\tabcolsep) * \real{0.5926}}
  >{\raggedright\arraybackslash}p{(\linewidth - 2\tabcolsep) * \real{0.4074}}@{}}
\toprule\noalign{}
\begin{minipage}[b]{\linewidth}\raggedright
Característica
\end{minipage} & \begin{minipage}[b]{\linewidth}\raggedright
Descrição
\end{minipage} \\
\midrule\noalign{}
\endhead
\bottomrule\noalign{}
\endlastfoot
\textbf{Fonte} & Kaggle - Data Hackers \\
\textbf{Período} & 2021 a 2024 \\
\textbf{Formato Original} & CSV (5 arquivos separados) \\
\textbf{Formato Final} & Parquet (consolidado) \\
\textbf{Total de Registros} & 17.419 respondentes \\
\textbf{Total de Colunas} & 1.225 \\
\textbf{Principais Variáveis} & UF, cargo, salário, área de atuação,
escolaridade, satisfação \\
\end{longtable}
}

\subsubsection{Variáveis Relevantes para o
Estudo}\label{variuxe1veis-relevantes-para-o-estudo}

\begin{itemize}
\tightlist
\item
  \textbf{UF} - Unidade Federativa onde o profissional reside
\item
  \textbf{Cargo} - Posição ocupada na organização
\item
  \textbf{Salário} - Faixa salarial do profissional
\item
  \textbf{Área de Atuação} - Setor da economia onde trabalha
\item
  \textbf{Escolaridade} - Nível de formação acadêmica
\item
  \textbf{Satisfação} - Nível de satisfação com o trabalho
\end{itemize}

\subsection{Base 2: Dados Socioeconômicos (API do
IBGE)}\label{base-2-dados-socioeconuxf4micos-api-do-ibge}

\subsubsection{Sobre a Fonte}\label{sobre-a-fonte}

O Instituto Brasileiro de Geografia e Estatística (IBGE) disponibiliza
uma API pública com dados oficiais sobre as Unidades Federativas
brasileiras, incluindo informações demográficas, econômicas e sociais.

\subsubsection{Características da
Base}\label{caracteruxedsticas-da-base-1}

{\def\LTcaptype{none} % do not increment counter
\begin{longtable}[]{@{}
  >{\raggedright\arraybackslash}p{(\linewidth - 2\tabcolsep) * \real{0.5926}}
  >{\raggedright\arraybackslash}p{(\linewidth - 2\tabcolsep) * \real{0.4074}}@{}}
\toprule\noalign{}
\begin{minipage}[b]{\linewidth}\raggedright
Característica
\end{minipage} & \begin{minipage}[b]{\linewidth}\raggedright
Descrição
\end{minipage} \\
\midrule\noalign{}
\endhead
\bottomrule\noalign{}
\endlastfoot
\textbf{Fonte} & API do IBGE - Localidades \\
\textbf{Endpoint} &
https://servicodados.ibge.gov.br/api/v1/localidades/estados \\
\textbf{Formato} & JSON \\
\textbf{Total de Registros} & 27 estados \\
\textbf{Total de Colunas} & 9 \\
\end{longtable}
}

\subsubsection{Variáveis Coletadas}\label{variuxe1veis-coletadas}

{\def\LTcaptype{none} % do not increment counter
\begin{longtable}[]{@{}
  >{\raggedright\arraybackslash}p{(\linewidth - 4\tabcolsep) * \real{0.3333}}
  >{\raggedright\arraybackslash}p{(\linewidth - 4\tabcolsep) * \real{0.3667}}
  >{\raggedright\arraybackslash}p{(\linewidth - 4\tabcolsep) * \real{0.3000}}@{}}
\toprule\noalign{}
\begin{minipage}[b]{\linewidth}\raggedright
Variável
\end{minipage} & \begin{minipage}[b]{\linewidth}\raggedright
Descrição
\end{minipage} & \begin{minipage}[b]{\linewidth}\raggedright
Período
\end{minipage} \\
\midrule\noalign{}
\endhead
\bottomrule\noalign{}
\endlastfoot
\textbf{UF} & Unidade Federativa (sigla) & 2024 \\
\textbf{População} & Estimativa populacional & 2024 \\
\textbf{PIB per capita} & Produto Interno Bruto por habitante & 2021 \\
\textbf{Renda Média} & Renda média da população & 2023 \\
\textbf{IDH} & Índice de Desenvolvimento Humano & 2021 \\
\textbf{Índice de Custo de Vida} & Estimativa de custo de vida relativo
& 2024 \\
\textbf{Poder de Compra} & Renda ajustada pelo custo de vida & 2024 \\
\textbf{Atratividade} & Medida de atratividade para profissionais &
2024 \\
\end{longtable}
}

\begin{quote}
\textbf{Nota:} Para variáveis como IDH e renda média, foram utilizados
dados de fontes secundárias (PNUD/IBGE) quando não disponíveis
diretamente pela API, seguindo a metodologia do trabalho.
\end{quote}

\begin{center}\rule{0.5\linewidth}{0.5pt}\end{center}

\section{Coleta de Dados}\label{coleta-de-dados}

\subsection{Coleta do State of Data
Brazil}\label{coleta-do-state-of-data-brazil}

\subsubsection{Método}\label{muxe9todo}

Os dados do State of Data foram coletados do Kaggle utilizando a
biblioteca \texttt{kagglehub}, que permite o download programático de
datasets.

\subsubsection{Script Utilizado}\label{script-utilizado}

```python \# src/coleta/consolida\_state\_of\_data.py

import kagglehub import pandas as pd from pathlib import Path

\section{Download automático do
dataset}\label{download-automuxe1tico-do-dataset}

caminho\_cache = kagglehub.dataset\_download(
``datahackers/state-of-data-brazil-20242025'' )

\section{Consolidação de todas as
edições}\label{consolidauxe7uxe3o-de-todas-as-ediuxe7uxf5es}

arquivos = sorted(Path(caminho\_cache).glob(``*.csv'')) quadros = {[}{]}
for arq in arquivos: df = pd.read\_csv(arq, low\_memory=False)
df{[}``arquivo\_origem''{]} = arq.name quadros.append(df)

df\_consolidado = pd.concat(quadros, ignore\_index=True)
df\_consolidado.to\_parquet(``data/processed/dados\_state\_of\_data\_consolidados.parquet'')
Desafios e Soluções Desafio Solução Arquivos em formatos diferentes
Padronização para CSV durante a leitura Colunas com nomes inconsistentes
Manutenção dos nomes originais com identificação do ano Dados com tipos
mistos Conversão para string antes do salvamento em Parquet Arquivo 2025
não disponível Utilização dos dados de 2024 como base mais recente
Coleta da Base Externa (API do IBGE) Método

A coleta foi realizada via requisições HTTP à API do IBGE, utilizando a
biblioteca requests do Python. Script Utilizado python

\section{src/coleta/coleta\_externa.py}\label{srccoletacoleta_externa.py}

import requests import pandas as pd

def coletar\_estados\_ibge(): url =
``https://servicodados.ibge.gov.br/api/v1/localidades/estados'' response
= requests.get(url, headers=HEADERS, timeout=30) dados = response.json()

\begin{verbatim}
estados = []
for estado in dados:
    estados.append({
        'uf': estado['sigla'],
        'estado': estado['nome'],
        'regiao': estado['regiao']['nome'],
        'populacao': estado.get('populacao', None)
    })

return pd.DataFrame(estados)
\end{verbatim}

Estrutura da Resposta da API json

\{ ``id'': 35, ``sigla'': ``SP'', ``nome'': ``São Paulo'', ``regiao'':
\{ ``id'': 3, ``nome'': ``Sudeste'' \}, ``populacao'': 45919049 \}

Tratamento de Dados Faltantes

Para variáveis não disponíveis diretamente na API, foram utilizadas
fontes complementares: Variável Fonte Complementar IDH PNUD/IBGE - Atlas
do Desenvolvimento Humano Renda Média PNAD/IBGE PIB per capita IBGE -
Contas Regionais Limpeza, Transformação e Integração dos Dados Limpeza
dos Dados State of Data

\begin{verbatim}
Remoção de Duplicatas: 7 duplicatas exatas foram removidas

Padronização de UF: Extração da sigla do formato "CEARÁ (CE)" para "CE"

Tratamento de Valores Nulos: Remoção de registros sem UF válida

Tipos de Dados: Conversão de colunas problemáticas para string
\end{verbatim}

Base Externa

\begin{verbatim}
Padronização de UF: Garantia de siglas em maiúsculas (SP, RJ, etc.)

Tratamento de Valores Nulos: Preenchimento com dados de fontes secundárias

Criação de Métricas Derivadas: Cálculo de poder de compra e atratividade
\end{verbatim}

Transformação dos Dados Criação de Métricas Derivadas Métrica Fórmula
Interpretação Poder de Compra Renda / Custo de Vida Poder aquisitivo
ajustado Atratividade Renda / (Custo de Vida)² Atratividade para
profissionais Profissionais per capita Profissionais / População *
100.000 Concentração relativa Padronização de UF

Um dos principais desafios foi a padronização dos nomes das Unidades
Federativas, que apareciam em formatos diferentes nas bases: python

\section{Função de extração de
sigla}\label{funuxe7uxe3o-de-extrauxe7uxe3o-de-sigla}

def extrair\_sigla(valor): \# ``CEARÁ (CE)'' -\textgreater{} ``CE''
match = re.search(r'(({[}A-Z{]}\{2\}))', valor) if match: return
match.group(1) return valor

Integração das Bases Estratégia de Integração

A integração foi realizada utilizando UF (Unidade Federativa) como chave
primária, através de um INNER JOIN. python

\section{src/limpeza/integra.py}\label{srclimpezaintegra.py}

\section{Agregar State of Data por
UF}\label{agregar-state-of-data-por-uf}

df\_sod\_agregado =
df\_sod.groupby(`uf').size().reset\_index(name=`total\_profissionais')

\section{Integrar com a base externa}\label{integrar-com-a-base-externa}

df\_integrado = pd.merge( df\_sod\_agregado, df\_ext, on=`uf',
how=`inner' )

Resultado da Integração Estado Profissionais População PIB per capita
Renda Média Custo de Vida IDH SP 1.401 45.919.049 R\$ 45.000 R\$ 3.500
1.45 0.835 MG 340 21.411.923 R\$ 34.000 R\$ 2.700 1.15 0.800 RJ 222
17.366.189 R\$ 42.000 R\$ 3.200 1.35 0.828 PR 149 11.597.484 R\$ 35.000
R\$ 2.800 1.08 0.812 DF 96 3.094.325 R\$ 48.000 R\$ 3.800 1.30 0.832

Total de estados integrados: 20 Modelo de Recuperação da Informação
Modelo Escolhido: Vetorial (TF-IDF + Similaridade de Cosseno)
Justificativa da Escolha

O modelo vetorial foi selecionado por:

\begin{verbatim}
Simplicidade e Interpretabilidade: Facilidade de explicar os resultados

Eficácia Comprovada: Ampla utilização em sistemas de busca

Adequação ao Problema: Busca por similaridade entre estados com base em características

Visualização Clara: Possibilidade de rankear resultados por similaridade
\end{verbatim}

Implementação python

from sklearn.feature\_extraction.text import TfidfVectorizer from
sklearn.metrics.pairwise import cosine\_similarity

\section{Criar documentos com características dos
estados}\label{criar-documentos-com-caracteruxedsticas-dos-estados}

df{[}`documento'{]} = df.apply(criar\_documento, axis=1)

\section{Criar matriz TF-IDF}\label{criar-matriz-tf-idf}

vectorizer = TfidfVectorizer(max\_features=100, stop\_words=`english')
tfidf\_matrix = vectorizer.fit\_transform(df{[}`documento'{]})

\section{Calcular similaridade de
cosseno}\label{calcular-similaridade-de-cosseno}

similaridades = cosine\_similarity(tfidf\_matrix{[}idx{]},
tfidf\_matrix)

Funcionamento do Modelo

\begin{verbatim}
Criação de Documentos: Cada estado é transformado em um documento textual com suas características

Cálculo do TF-IDF: Ponderar a importância de cada termo no documento e na coleção

Similaridade de Cosseno: Medir a distância entre o vetor do estado de referência e os demais

Ranking: Ordenar os estados por similaridade decrescente
\end{verbatim}

Exemplo de Documento Criado text

Estado: SP Região: Sudeste Profissionais: 1401 Renda média: R\$ 3500
Custo de vida: 1.45 PIB per capita: R\$ 45000 IDH: 0.835 Poder de
compra: 2413.79

Resultados da Busca Estado Referência Estados Mais Similares
Similaridade SP MG (0.291), ES (0.267), MA (0.247) Padrão Sudeste DF GO
(0.333), MS (0.308), PR (0.300) Centro-Oeste BA PE (0.429), GO (0.313),
MA (0.302) Nordeste Termos Mais Relevantes (TF-IDF) Termo TF-IDF Médio
Importância profissionais 0.4234 Diferencia estados por número de
profissionais renda 0.3987 Diferencia por nível econômico pib 0.3654
Diferencia por desenvolvimento custo 0.3421 Diferencia por custo de vida
idh 0.3218 Diferencia por desenvolvimento humano Visualizações e Análise
dos Resultados 1. Distribuição de Profissionais por Estado

https://docs/figuras/top10\_profissionais.png\{\#fig-top10 width=100\%\}
Análise

\begin{verbatim}
Concentração Regional: Os estados do Sudeste (SP, MG, RJ) concentram a maior parte dos profissionais

Destaque para SP: São Paulo possui mais profissionais que todos os outros estados combinados

Desigualdade Regional: A distribuição é extremamente desigual, com forte concentração no eixo Sul-Sudeste
\end{verbatim}

\begin{enumerate}
\def\labelenumi{\arabic{enumi}.}
\setcounter{enumi}{1}
\tightlist
\item
  Relação entre Profissionais e Custo de Vida
\end{enumerate}

https://docs/figuras/profissionais\_vs\_custo.png\{\#fig-custo
width=100\%\} Análise

\begin{verbatim}
Correlação Positiva: Estados com maior custo de vida tendem a ter mais profissionais

SP como Outlier: São Paulo destaca-se com alta concentração de profissionais

Relação com Renda: A renda média também parece estar correlacionada com a concentração
\end{verbatim}

\begin{enumerate}
\def\labelenumi{\arabic{enumi}.}
\setcounter{enumi}{2}
\tightlist
\item
  Matriz de Correlação
\end{enumerate}

https://docs/figuras/matriz\_correlacao.png\{\#fig-correlacao
width=100\%\} Correlações com Número de Profissionais Indicador
Correlação Interpretação Custo de Vida +0.65 Forte correlação positiva
PIB per capita +0.56 Correlação positiva moderada Renda Média +0.56
Correlação positiva moderada IDH +0.51 Correlação positiva moderada
Atratividade -0.61 Correlação negativa SURPREENDENTE Análise da
Correlação Negativa com Atratividade

A correlação negativa entre atratividade (renda/custo²) e número de
profissionais é um achado contra-intuitivo que pode ser explicado por:

\begin{verbatim}
Saturação do Mercado: Estados com maior atratividade já estão saturados de profissionais

Efeito de Aglomeração: Profissionais migram para centros já estabelecidos, mesmo com custo mais alto

Melhores Oportunidades: Estados com maior custo de vida oferecem melhores salários, atraindo profissionais

Infraestrutura: Estados mais desenvolvidos oferecem melhores condições de trabalho
\end{verbatim}

\begin{enumerate}
\def\labelenumi{\arabic{enumi}.}
\setcounter{enumi}{3}
\tightlist
\item
  Distribuição por Região
\end{enumerate}

https://docs/figuras/profissionais\_por\_regiao.png\{\#fig-regiao
width=100\%\} Análise

\begin{verbatim}
Concentração no Sudeste: A região Sudeste concentra a maioria absoluta dos profissionais

Centro-Oeste em Destaque: Apesar de menor população, tem presença significativa de profissionais

Sub-representação do Norte: A região Norte tem a menor concentração de profissionais
\end{verbatim}

\begin{enumerate}
\def\labelenumi{\arabic{enumi}.}
\setcounter{enumi}{4}
\tightlist
\item
  Relação entre PIB e Profissionais
\end{enumerate}

https://docs/figuras/pib\_vs\_profissionais.png\{\#fig-pib width=100\%\}
Análise

\begin{verbatim}
Relação Positiva: Estados com maior PIB per capita têm mais profissionais

SP como Destaque: São Paulo é líder em ambos os indicadores

Relação com Custo de Vida: Estados com maior PIB também tendem a ter maior custo de vida
\end{verbatim}

\begin{enumerate}
\def\labelenumi{\arabic{enumi}.}
\setcounter{enumi}{5}
\tightlist
\item
  Ranking de Atratividade
\end{enumerate}

https://docs/figuras/ranking\_atratividade.png\{\#fig-atratividade
width=100\%\} Análise

\begin{verbatim}
DF em Destaque: Distrito Federal tem a maior atratividade

Sudeste Líder: SP, RJ e MG estão entre os mais atrativos

Norte com Menor Atratividade: AM e PA têm a menor atratividade
\end{verbatim}

Conclusões Principais Achados 1. Concentração Regional Significativa

A pesquisa revelou uma forte concentração de profissionais de dados na
região Sudeste, especialmente no estado de São Paulo. Essa concentração
está alinhada com a distribuição do PIB e da atividade econômica no
país. 2. Fatores Econômicos Determinantes

Os principais fatores associados à concentração de profissionais são:

\begin{verbatim}
PIB per capita (correlação 0.56)

Renda Média (correlação 0.56)

Custo de Vida (correlação 0.65)

IDH (correlação 0.51)
\end{verbatim}

\begin{enumerate}
\def\labelenumi{\arabic{enumi}.}
\setcounter{enumi}{2}
\tightlist
\item
  Efeito Surpresa da Atratividade
\end{enumerate}

A correlação negativa entre atratividade (renda/custo²) e número de
profissionais (-0.61) sugere que:

\begin{verbatim}
Estados com maior atratividade podem estar saturados

Profissionais buscam oportunidades mesmo com custo de vida mais alto

O efeito de aglomeração supera o custo-benefício individual
\end{verbatim}

\begin{enumerate}
\def\labelenumi{\arabic{enumi}.}
\setcounter{enumi}{3}
\tightlist
\item
  Modelo TF-IDF Eficaz
\end{enumerate}

O modelo de RI aplicado conseguiu identificar estados com perfis
similares, agrupando corretamente estados da mesma região e com
características econômicas semelhantes. Contribuições do Trabalho

\begin{verbatim}
Metodologia Reproduzível: Todo o processo de coleta e análise é documentado e automatizado

Base de Dados Integrada: Criação de um dataset unificado com 20 estados e 13 variáveis

Aplicação de RI: Demonstração prática do modelo TF-IDF em dados reais

Insights Relevantes: Identificação de padrões da distribuição geográfica
\end{verbatim}

Limitações Limitações dos Dados 1. Cobertura Incompleta

\begin{verbatim}
Apenas 20 estados integrados: 7 estados não tinham registros no State of Data

Ausência do ano de 2025: A edição mais recente não estava disponível

Dados de custo de vida estimados: Não são dados oficiais do IBGE
\end{verbatim}

\begin{enumerate}
\def\labelenumi{\arabic{enumi}.}
\setcounter{enumi}{1}
\item
  Viés da Pesquisa State of Data

  Auto-seleção: Respondentes são voluntários, podendo haver viés de
  amostragem

  Representatividade: A amostra pode não ser completamente
  representativa

  Perguntas em inglês: Algumas questões em inglês podem ter influenciado
  respostas
\item
  Limitações da Base Externa

  Dados agregados: Não há granularidade por cidade ou região

  Períodos diferentes: As variáveis não são do mesmo ano

  Dados faltantes: Algumas variáveis foram preenchidas com estimativas
\end{enumerate}

Limitações Metodológicas 1. Modelo TF-IDF Simplificado

\begin{verbatim}
Sem contexto semântico: O modelo não captura significado, apenas frequência

Stop words removidas: Pode ter perdido informações importantes

Poucas variáveis: Apenas 7 características foram consideradas
\end{verbatim}

\begin{enumerate}
\def\labelenumi{\arabic{enumi}.}
\setcounter{enumi}{1}
\item
  Análise Estatística Básica

  Correlação não é causalidade: As correlações não indicam relação de
  causa e efeito

  Dados agregados: Perde-se a heterogeneidade dentro dos estados

  Sem testes de significância: Não foram aplicados testes estatísticos
  formais
\end{enumerate}

Limitações Técnicas 1. Infraestrutura

\begin{verbatim}
Capacidade computacional: Processamento de 1.225 colunas exigiu otimizações

Memória RAM: Em alguns momentos foi necessário reduzir o dataset
\end{verbatim}

\begin{enumerate}
\def\labelenumi{\arabic{enumi}.}
\setcounter{enumi}{1}
\item
  Reproducibilidade

  Dependência externa: O Kaggle pode mudar a estrutura dos dados

  API do IBGE: Pode sofrer alterações ou ficar indisponível
\end{enumerate}

Trabalhos Futuros

\begin{verbatim}
Inclusão de mais variáveis: Educação, infraestrutura, conectividade

Modelos mais avançados: BM25, word embeddings, BERT

Análise por cidade: Maior granularidade

Série temporal: Análise da evolução ao longo dos anos

Dashboard interativo: Visualização dinâmica dos dados
\end{verbatim}

Referências

\begin{verbatim}
Data Hackers. (2021-2024). State of Data Brazil. Kaggle. Disponível em: https://www.kaggle.com/datasets/datahackers

IBGE - Instituto Brasileiro de Geografia e Estatística. (2024). API de Localidades. Disponível em: https://servicodados.ibge.gov.br/api/docs/localidades

IBGE. (2021). Produto Interno Bruto dos Municípios. Disponível em: https://www.ibge.gov.br/estatisticas/economicas/contas-nacionais

PNUD. (2021). Atlas do Desenvolvimento Humano no Brasil. Disponível em: http://www.atlasbrasil.org.br/

PNAD - Pesquisa Nacional por Amostra de Domicílios. (2023). IBGE. Disponível em: https://www.ibge.gov.br/estatisticas/sociais/trabalho

Manning, C. D., Raghavan, P., & Schütze, H. (2008). Introduction to Information Retrieval. Cambridge University Press.

Baeza-Yates, R., & Ribeiro-Neto, B. (2011). Modern Information Retrieval: The Concepts and Technology Behind Search. Addison-Wesley.
\end{verbatim}

Apêndice A. Scripts Utilizados A.1 Consolidação do State of Data python

\section{src/coleta/consolida\_state\_of\_data.py}\label{srccoletaconsolida_state_of_data.py}

\section{Código completo disponível no
repositório}\label{cuxf3digo-completo-disponuxedvel-no-reposituxf3rio}

A.2 Coleta da Base Externa python

\section{src/coleta/coleta\_externa.py}\label{srccoletacoleta_externa.py-1}

\section{Código completo disponível no
repositório}\label{cuxf3digo-completo-disponuxedvel-no-reposituxf3rio-1}

A.3 Integração das Bases python

\section{src/limpeza/integra.py}\label{srclimpezaintegra.py-1}

\section{Código completo disponível no
repositório}\label{cuxf3digo-completo-disponuxedvel-no-reposituxf3rio-2}

A.4 Modelo de RI (TF-IDF) python

\section{modelo\_ri.py}\label{modelo_ri.py}

\section{Código completo disponível no
repositório}\label{cuxf3digo-completo-disponuxedvel-no-reposituxf3rio-3}

B. Dados Gerados B.1 Base Integrada

\begin{verbatim}
Arquivo: data/processed/base_integrada.parquet

20 estados x 13 colunas
\end{verbatim}

B.2 Visualizações

\begin{verbatim}
docs/figuras/top10_profissionais.png

docs/figuras/profissionais_vs_custo.png

docs/figuras/matriz_correlacao.png

docs/figuras/profissionais_por_regiao.png

docs/figuras/pib_vs_profissionais.png

docs/figuras/ranking_atratividade.png
\end{verbatim}

B.3 Resultados do Modelo RI

\begin{verbatim}
docs/resultados_busca.csv
\end{verbatim}




\end{document}
